\chapter{Implementing ROS2 with Python}

\begin{lernziele}
After this chapter, you understand the structure of a simple ROS2 Python node. You can place publishers, subscribers, timers, and basic package structures, and you know how to derive a simple node from the model.
\end{lernziele}

\section{Introduction}
Once the architecture is modeled, implementation follows. In this chapter, we look at Python with \texttt{rclpy}. Python suits beginners well, because its code is comparatively easy to read.

\section{Creating a Package}
You typically create a Python package with the following command:

\begin{lstlisting}[language=bash]
ros2 pkg create robot_basic_nodes --build-type ament_python --dependencies rclpy std_msgs sensor_msgs geometry_msgs
\end{lstlisting}

The parameter \texttt{--build-type ament\_python} sets the build type: \texttt{ament\_python} suits pure Python packages, while \texttt{ament\_cmake} suits C++ or mixed packages.

A package bundles code, configuration, and startup information. Larger systems benefit from several packages. Figure~\ref{fig:kap17-workspace-struktur} shows the resulting workspace and package structure.

\begin{figure}[H]
  \centering
  \resizebox{\textwidth}{!}{% =====================================================================
% Figure: ROS2 workspace and package structure
% Matches Chapter 17, section "Creating a Package"
%
% Included in the book, e.g. in chapters_en/17_ros2_python.tex:
%
%   \begin{figure}[H]
%     \centering
%     \resizebox{\textwidth}{!}{% =====================================================================
% Figure: ROS2 workspace and package structure
% Matches Chapter 17, section "Creating a Package"
%
% Included in the book, e.g. in chapters_en/17_ros2_python.tex:
%
%   \begin{figure}[H]
%     \centering
%     \resizebox{\textwidth}{!}{% =====================================================================
% Figure: ROS2 workspace and package structure
% Matches Chapter 17, section "Creating a Package"
%
% Included in the book, e.g. in chapters_en/17_ros2_python.tex:
%
%   \begin{figure}[H]
%     \centering
%     \resizebox{\textwidth}{!}{\input{figures_en/fig_kap17_workspace_struktur}}
%     \caption{Directory structure of a ROS2 workspace with the package
%       \texttt{robot\_basic\_nodes} and its sibling folders
%       \texttt{build/}, \texttt{install/}, and \texttt{log/}.}
%     \label{fig:kap17-workspace-struktur}
%   \end{figure}
%
% Required packages (already loaded in the book via main_en.tex): tikz, graphicx
% Required TikZ libraries: positioning, arrows.meta, shadows, calc
% =====================================================================

\input{figures_en/fig_kap17_workspace_struktur.tikz}
}
%     \caption{Directory structure of a ROS2 workspace with the package
%       \texttt{robot\_basic\_nodes} and its sibling folders
%       \texttt{build/}, \texttt{install/}, and \texttt{log/}.}
%     \label{fig:kap17-workspace-struktur}
%   \end{figure}
%
% Required packages (already loaded in the book via main_en.tex): tikz, graphicx
% Required TikZ libraries: positioning, arrows.meta, shadows, calc
% =====================================================================

\begin{tikzpicture}[
    scale=0.85, transform shape,
    every node/.style={font=\small},
    folder/.style={
        rectangle, rounded corners=1.5pt, draw=black!70, thick,
        minimum width=34mm, minimum height=9mm,
        align=center, fill=blue!8, drop shadow
    },
    subfolder/.style={folder, minimum width=44mm, fill=green!12, draw=green!40!black},
    filenode/.style={folder, minimum width=48mm, minimum height=8mm, fill=orange!12,
        draw=orange!60!black, font=\small},
    sibling/.style={folder, fill=black!6, draw=black!50},
    trunkline/.style={thick, black!55},
    flow/.style={-{Stealth[length=2mm]}, thick, black!60},
    lbl/.style={font=\small}
]

% --- Root ---
\node[folder] (ws) at (0, 0) {\texttt{robot\_ws/}};

% --- Level 1 (compact column, tree connection via a side trunk) ---
\node[subfolder] (src)     at (5.0, 1.95)  {\texttt{src/robot\_basic\_nodes/}};
\node[sibling]   (build)   at (5.0, 0.65)  {\texttt{build/}};
\node[sibling]   (install) at (5.0, -0.65) {\texttt{install/}};
\node[sibling]   (log)     at (5.0, -1.95) {\texttt{log/}};

\draw[trunkline] (ws.east) -- (2.2,0 |- ws.east);
\draw[trunkline] (2.2,0 |- src.west) -- (2.2,0 |- log.west);
\draw[flow] (2.2,0 |- src.west) -- (src.west);
\draw[flow] (2.2,0 |- build.west) -- (build.west);
\draw[flow] (2.2,0 |- install.west) -- (install.west);
\draw[flow] (2.2,0 |- log.west) -- (log.west);

% --- Level 2: contents of src/robot_basic_nodes (its own trunk) ---
\node[filenode] (setup)    at (10.6, 1.3)  {\texttt{setup.py}};
\node[filenode] (pkgxml)   at (10.6, 0)    {\texttt{package.xml}};
\node[filenode] (statuspy) at (10.6, -1.3) {\texttt{robot\_basic\_nodes/}\\\texttt{status\_node.py}};

\draw[trunkline] (src.east) -- (7.7,0 |- src.east);
\draw[trunkline] (7.7,0 |- setup.west) -- (7.7,0 |- statuspy.west);
\draw[flow] (7.7,0 |- setup.west) -- (setup.west);
\draw[flow] (7.7,0 |- pkgxml.west) -- (pkgxml.west);
\draw[flow] (7.7,0 |- statuspy.west) -- (statuspy.west);

% --- Note on install/setup.bash ---
\node[lbl, align=left, below=10mm of log, xshift=18mm, text width=48mm]
    (note) {loaded with \texttt{source}, so that \texttt{ros2 run} finds the package};
\draw[flow, dashed] (install.south) -- ++(0,-0.6) -| (note.north);

\end{tikzpicture}

}
%     \caption{Directory structure of a ROS2 workspace with the package
%       \texttt{robot\_basic\_nodes} and its sibling folders
%       \texttt{build/}, \texttt{install/}, and \texttt{log/}.}
%     \label{fig:kap17-workspace-struktur}
%   \end{figure}
%
% Required packages (already loaded in the book via main_en.tex): tikz, graphicx
% Required TikZ libraries: positioning, arrows.meta, shadows, calc
% =====================================================================

\begin{tikzpicture}[
    scale=0.85, transform shape,
    every node/.style={font=\small},
    folder/.style={
        rectangle, rounded corners=1.5pt, draw=black!70, thick,
        minimum width=34mm, minimum height=9mm,
        align=center, fill=blue!8, drop shadow
    },
    subfolder/.style={folder, minimum width=44mm, fill=green!12, draw=green!40!black},
    filenode/.style={folder, minimum width=48mm, minimum height=8mm, fill=orange!12,
        draw=orange!60!black, font=\small},
    sibling/.style={folder, fill=black!6, draw=black!50},
    trunkline/.style={thick, black!55},
    flow/.style={-{Stealth[length=2mm]}, thick, black!60},
    lbl/.style={font=\small}
]

% --- Root ---
\node[folder] (ws) at (0, 0) {\texttt{robot\_ws/}};

% --- Level 1 (compact column, tree connection via a side trunk) ---
\node[subfolder] (src)     at (5.0, 1.95)  {\texttt{src/robot\_basic\_nodes/}};
\node[sibling]   (build)   at (5.0, 0.65)  {\texttt{build/}};
\node[sibling]   (install) at (5.0, -0.65) {\texttt{install/}};
\node[sibling]   (log)     at (5.0, -1.95) {\texttt{log/}};

\draw[trunkline] (ws.east) -- (2.2,0 |- ws.east);
\draw[trunkline] (2.2,0 |- src.west) -- (2.2,0 |- log.west);
\draw[flow] (2.2,0 |- src.west) -- (src.west);
\draw[flow] (2.2,0 |- build.west) -- (build.west);
\draw[flow] (2.2,0 |- install.west) -- (install.west);
\draw[flow] (2.2,0 |- log.west) -- (log.west);

% --- Level 2: contents of src/robot_basic_nodes (its own trunk) ---
\node[filenode] (setup)    at (10.6, 1.3)  {\texttt{setup.py}};
\node[filenode] (pkgxml)   at (10.6, 0)    {\texttt{package.xml}};
\node[filenode] (statuspy) at (10.6, -1.3) {\texttt{robot\_basic\_nodes/}\\\texttt{status\_node.py}};

\draw[trunkline] (src.east) -- (7.7,0 |- src.east);
\draw[trunkline] (7.7,0 |- setup.west) -- (7.7,0 |- statuspy.west);
\draw[flow] (7.7,0 |- setup.west) -- (setup.west);
\draw[flow] (7.7,0 |- pkgxml.west) -- (pkgxml.west);
\draw[flow] (7.7,0 |- statuspy.west) -- (statuspy.west);

% --- Note on install/setup.bash ---
\node[lbl, align=left, below=10mm of log, xshift=18mm, text width=48mm]
    (note) {loaded with \texttt{source}, so that \texttt{ros2 run} finds the package};
\draw[flow, dashed] (install.south) -- ++(0,-0.6) -| (note.north);

\end{tikzpicture}

}
  \caption{Directory structure of a ROS2 workspace with the package \texttt{robot\_basic\_nodes} and its sibling folders \texttt{build/}, \texttt{install/}, and \texttt{log/}.}
  \label{fig:kap17-workspace-struktur}
\end{figure}

\section{A Simple Node}
\begin{lstlisting}[language=Python]
import rclpy
from rclpy.node import Node
from std_msgs.msg import String

class StatusNode(Node):
    def __init__(self):
        super().__init__('status_node')
        self.publisher = self.create_publisher(String, '/robot/status', 10)
        self.timer = self.create_timer(1.0, self.publish_status)

    def publish_status(self):
        msg = String()
        msg.data = 'Robot is running'
        self.publisher.publish(msg)

def main():
    rclpy.init()
    node = StatusNode()
    rclpy.spin(node)
    node.destroy_node()
    rclpy.shutdown()

if __name__ == '__main__':
    main()
\end{lstlisting}

For \texttt{ros2 run robot\_basic\_nodes status\_node} to find this code later, you must also add a matching entry to the \texttt{entry\_points} of the package's \texttt{setup.py}; we do not go into that in detail here.

\section{Publisher}
A publisher publishes messages on a topic. In the example, the node regularly sends a status on \texttt{/robot/status}. A battery monitor would publish \texttt{/battery\_state} in a similarly regular way.

\section{Subscriber}
A subscriber receives messages from a topic. The ObjectDetector, for example, would subscribe to \texttt{/camera/image\_raw}. A callback function (a method that is called automatically as soon as a new message arrives) processes incoming messages.

\section{Timer}
Timers trigger functions on a regular schedule. They suit status messages, simulated sensor data, or periodic checks. Figure~\ref{fig:kap17-ablauf-publisher-timer} shows the sequence inside \texttt{StatusNode}, from the timer to the message on the topic.

\begin{figure}[H]
  \centering
  \adjustbox{max width=\linewidth,max totalheight=0.6\textheight}{% =====================================================================
% Figure: Publisher/timer sequence in StatusNode
% Matches Chapter 17, section "Timer"
%
% Included in the book, e.g. in chapters_en/17_ros2_python.tex:
%
%   \begin{figure}[H]
%     \centering
%     \resizebox{\textwidth}{!}{% =====================================================================
% Figure: Publisher/timer sequence in StatusNode
% Matches Chapter 17, section "Timer"
%
% Included in the book, e.g. in chapters_en/17_ros2_python.tex:
%
%   \begin{figure}[H]
%     \centering
%     \resizebox{\textwidth}{!}{% =====================================================================
% Figure: Publisher/timer sequence in StatusNode
% Matches Chapter 17, section "Timer"
%
% Included in the book, e.g. in chapters_en/17_ros2_python.tex:
%
%   \begin{figure}[H]
%     \centering
%     \resizebox{\textwidth}{!}{\input{figures_en/fig_kap17_ablauf_publisher_timer}}
%     \caption{Sequence inside \texttt{StatusNode}: the timer periodically
%       triggers \texttt{publish\_status()}, which creates a message and
%       publishes it on the topic \texttt{/robot/status}.}
%     \label{fig:kap17-ablauf-publisher-timer}
%   \end{figure}
%
% Required packages (already loaded in the book via main_en.tex): tikz, graphicx
% Required TikZ libraries: positioning, arrows.meta, shadows, calc
% =====================================================================

\input{figures_en/fig_kap17_ablauf_publisher_timer.tikz}
}
%     \caption{Sequence inside \texttt{StatusNode}: the timer periodically
%       triggers \texttt{publish\_status()}, which creates a message and
%       publishes it on the topic \texttt{/robot/status}.}
%     \label{fig:kap17-ablauf-publisher-timer}
%   \end{figure}
%
% Required packages (already loaded in the book via main_en.tex): tikz, graphicx
% Required TikZ libraries: positioning, arrows.meta, shadows, calc
% =====================================================================

\begin{tikzpicture}[
    every node/.style={font=\small},
    ros/.style={
        rectangle, rounded corners=1.5pt, draw=orange!60!black, thick,
        minimum width=42mm, minimum height=11mm,
        align=center, fill=orange!15, drop shadow
    },
    ext/.style={
        rectangle, rounded corners=1.5pt, draw=black!50, thick, dashed,
        minimum width=42mm, minimum height=13mm,
        align=center, fill=white, drop shadow
    },
    flow/.style={-{Stealth[length=2.2mm]}, thick, black!60},
    dflow/.style={-{Stealth[length=2.2mm]}, thick, orange!70!black, dashed},
    lbl/.style={font=\small, fill=white, inner sep=1pt}
]

\node[ros] (timer) at (0, 0)     {Timer\\(period 1.0 s)};
\node[ros] (pub)   at (0, -2.0)  {\texttt{publish\_status()}};
\node[ros] (msg)   at (0, -4.0)  {String message\\\texttt{'Robot is running'}};
\node[ros] (topic) at (0, -6.0)  {Topic\\\texttt{/robot/status}};
\node[ext] (sub)   at (0, -8.3)  {Subscriber node\\(e.g.\ \texttt{user\_interface\_}\\\texttt{node})};

\draw[flow] (timer.south) -- node[lbl, right]{triggers} (pub.north);
\draw[flow] (pub.south)   -- node[lbl, right]{creates}  (msg.north);
\draw[flow] (msg.south)   -- node[lbl, right]{\texttt{publisher.publish()}} (topic.north);
\draw[dflow] (topic.south) -- (sub.north);

\end{tikzpicture}

}
%     \caption{Sequence inside \texttt{StatusNode}: the timer periodically
%       triggers \texttt{publish\_status()}, which creates a message and
%       publishes it on the topic \texttt{/robot/status}.}
%     \label{fig:kap17-ablauf-publisher-timer}
%   \end{figure}
%
% Required packages (already loaded in the book via main_en.tex): tikz, graphicx
% Required TikZ libraries: positioning, arrows.meta, shadows, calc
% =====================================================================

\begin{tikzpicture}[
    every node/.style={font=\small},
    ros/.style={
        rectangle, rounded corners=1.5pt, draw=orange!60!black, thick,
        minimum width=42mm, minimum height=11mm,
        align=center, fill=orange!15, drop shadow
    },
    ext/.style={
        rectangle, rounded corners=1.5pt, draw=black!50, thick, dashed,
        minimum width=42mm, minimum height=13mm,
        align=center, fill=white, drop shadow
    },
    flow/.style={-{Stealth[length=2.2mm]}, thick, black!60},
    dflow/.style={-{Stealth[length=2.2mm]}, thick, orange!70!black, dashed},
    lbl/.style={font=\small, fill=white, inner sep=1pt}
]

\node[ros] (timer) at (0, 0)     {Timer\\(period 1.0 s)};
\node[ros] (pub)   at (0, -2.0)  {\texttt{publish\_status()}};
\node[ros] (msg)   at (0, -4.0)  {String message\\\texttt{'Robot is running'}};
\node[ros] (topic) at (0, -6.0)  {Topic\\\texttt{/robot/status}};
\node[ext] (sub)   at (0, -8.3)  {Subscriber node\\(e.g.\ \texttt{user\_interface\_}\\\texttt{node})};

\draw[flow] (timer.south) -- node[lbl, right]{triggers} (pub.north);
\draw[flow] (pub.south)   -- node[lbl, right]{creates}  (msg.north);
\draw[flow] (msg.south)   -- node[lbl, right]{\texttt{publisher.publish()}} (topic.north);
\draw[dflow] (topic.south) -- (sub.north);

\end{tikzpicture}

}
  \caption{Sequence inside \texttt{StatusNode}: the timer periodically triggers \texttt{publish\_status()}, which creates a message and publishes it on the topic \texttt{/robot/status}.}
  \label{fig:kap17-ablauf-publisher-timer}
\end{figure}

\section{Build and Start}
\texttt{colcon} is the standard build tool for ROS2 workspaces; it compiles and links all packages in the workspace together. The workspace typically splits into the folders \texttt{src/} (source code), \texttt{build/} and \texttt{install/} (build results), and \texttt{log/} (logs), as shown in Figure~\ref{fig:kap17-workspace-struktur}.

\begin{lstlisting}[language=bash]
colcon build
source install/setup.bash
ros2 run robot_basic_nodes status_node
\end{lstlisting}

\section{From Model to Node}
The \texttt{part def BatterySystem} can lead to the decision to implement a \texttt{battery\_monitor\_node}. The interface \enquote{battery status} becomes the topic \texttt{/battery\_state}. The test case later checks whether the node publishes plausible values.

\begin{hinweisbox}
The code here is deliberately minimal. Real projects add error handling, parameters, logging, launch files, tests, and a clean package structure.
\end{hinweisbox}

\section{Common Mistakes}
\subsection{No Clear Responsibility}
A node should have an understandable task.

\subsection{Hard-Coding Topic Names}
Topic names should stay consistent, and be configurable through parameters where needed.

\subsection{No Error Handling}
Even simple nodes should not ignore invalid data and exceptions.

\section{From a Minimal Example to Project Code}
The example shown here is deliberately small. In a real project, a node should do more: read parameters, use logging, detect invalid data, and shut down cleanly. For \texttt{battery\_monitor\_node}, parameters could define the percentage at which the node publishes a warning.

A sensible next step is a simulated battery node. It regularly publishes a battery level that slowly drops. A second node can subscribe to this value, and issue a warning when the level is low. This produces a small, testable ROS2 system that connects directly to requirements and model relationships.

The link back to the model matters. Code should not grow in isolation. Every node should have a purpose in the model, belong to a function, and communicate over defined topics.

\section{Summary}
ROS2 Python nodes implement parts of the modeled architecture. Publishers, subscribers, and timers form the basic building blocks. For a start, a small node is enough, as long as its role in the model is clear.

\section{Review Questions}
\begin{enumerate}
    \item What is a ROS2 package?
    \item What does a publisher do?
    \item What does a subscriber do?
    \item What is a timer for?
    \item Why should a node have a clear responsibility?
    \item How can a SysML system element (\texttt{part def}) lead to a ROS2 node?
\end{enumerate}

\section{Exercises}
\begin{uebungbox}
\textbf{Exercise 1: Battery Monitor}\\
Describe conceptually how you would implement a \texttt{battery\_monitor\_node} that regularly publishes a simulated battery level.
\end{uebungbox}

\begin{uebungbox}
\textbf{Exercise 2: Model Link}\\
Describe which requirement, which system element (\texttt{part def}), which topic, and which test case belong to your \texttt{battery\_monitor\_node}.
\end{uebungbox}

\section{Further Reading}
For practical ROS2 implementation, the official documentation \cite{ros2docs} is central. This chapter stays a deliberate starting point, and prepares you for the full ROS2 project later.
